%\appendix

\section{\gls{acr-meem}}
\label{sec:appendix-meem-details}
As introduced in \Cref{sec:hydro-meem}, the hydrodynamic coefficients are computed semi-analytically using \gls{acr-meem}.
This section explains the \gls{acr-meem} formulation and solution methodology for the radiation of a dual truncated concentric cylinder geometry, originally presented by \citet{mavrakos_hydrodynamic_2004,chau_inertia_2010,chau_inertia_2012} as an extension of the single-cylinder \gls{acr-meem} radiation solution in the study by \citet{yeung_added_1981}.
Further numeric and realization details of the authors' implementation may be found in references \citet{bimali_matrix_2026,best_openflash_2026,mccabe_open-source_2024}.
The computation involves splitting the fluid domain into regions, approximating an infinite series by truncation, and solving a matrix equation to enforce the continuity of potential and velocity across regions.

\subsection{Linear Hydrodynamics and Eigenfunctions}
The dynamics of a floating body in water waves are well-described by linear potential flow theory, a simplification of the Navier-Stokes equation.
This theory states that the fluid velocity field is the gradient of some complex potential $\gls{sym-phi}$, $\gls{sym-vec-v}=\nabla\gls{sym-phi}$, and $\gls{sym-phi}$ satisfies the Laplace equation, $\nabla^2\gls{sym-phi}=0$.
Adding the free surface condition, far-field or incident waves, and body surface and sea-bed conditions detailed in \citet{chatjigeorgiou_analytical_2018} yields a boundary value problem.
When boundary conditions correspond to the heave radiation problem (body moving vertically, no incident waves), solving for the potential $\gls{sym-phi}(\gls{sym-r},\gls{sym-theta},\gls{sym-z})$ determines the heave added mass and damping $\gls{sym-A-h}$ and $\gls{sym-B-h}$, hereafter ``hydro coefficients.”

For appropriate geometries, the partial differential equation is separable and $\gls{sym-phi}$ can be expressed as the product of radial, vertical, and circumferential basis functions called eigenfunctions.
In cylindrically symmetric problems, the radial eigenfunctions are a family of transcendental functions called Bessel functions.
The fluid is then divided into cylindrical regions.
Arbitrarily many fluid regions can exist, so the method applies to any axisymmetric geometry, including multiple concentric bodies that oscillate independently.
Here, two concentric cylinders and thus three fluid regions are demonstrated.
Extension to many regions is discussed in section 3.4.
\Cref{fig:meem-regions} illustrates the regions and dimensions: two interior regions \textit{i1} and \textit{i2}, and an exterior region \textit{e} extending to infinity.

\begin{figure}[htb]
\centering
\includegraphics[\FigureOptions{aor:figs/from-matlab/meem_dims.pdf}]{figs/from-matlab/meem_dims.pdf}
\caption{Fluid regions and dimensions used in the dual concentric cylinder \gls{acr-meem}}
\label{fig:meem-regions}
\end{figure}

 The potential in each region is split into a homogeneous part for the unforced solution and a particular part due to body motion: $\gls{sym-phi}=\gls{sym-phi-h}+\gls{sym-phi-p}$.
Boundary conditions dictate $\gls{sym-phi-p}$ and the eigenfunctions for $\gls{sym-phi-h}$ in each region, which the textbook \citet{chatjigeorgiou_analytical_2018} describes in detail.
\Cref{tab:MEEM-eigenfunctions} shows the equations originally presented in the studies by \citet{chau_inertia_2010,chau_inertia_2012} for the potential and eigenfunctions in each region, which include infinitely many unknown eigencoefficients $\gls{sym-C-1n-i1}$, $\gls{sym-C-1m-i2}$, $\gls{sym-C-2m-i2}$ and $\gls{sym-B-k-e}$.
By construction, this potential obeys all boundary conditions except for zero radial velocity on radial body surfaces.
The unknown coefficients must be computed to enforce this final condition as well as continuity across regions, which will be the subject of \Cref{sec:appendix-meem-matching}.

\newcommand{\dispersion}{
    \begin{cases}
        \gls{sym-m-0}\tanh (\gls{sym-m-0}\gls{sym-h})
            =\gls{sym-omega}^{2}/\gls{sym-g},
            & \gls{sym-k-index}=0,\\
        \gls{sym-m-k}\tan (\gls{sym-m-k}\gls{sym-h})
            =-\gls{sym-omega}^{2}/\gls{sym-g},
            & \gls{sym-k-index}\geq 1
    \end{cases}
}

\begin{table}[htbp]
    \centering
    \expanded{
        \noexpand\tabular{\TableColumnSpecs{aor:meem-eigenfunctions}}
    }
    \hline 
    Region
        &
        $i1$
        &
        $i2$
        &
        $e$
        \\
        \hline 
    Homog. potential $\gls{sym-phi-h}(\gls{sym-r},\gls{sym-z})$
        &
        $\displaystyle\sum_{\gls{sym-n}}
            \gls{sym-C-1n-i1}
            \gls{sym-R-1n-i1} (\gls{sym-r})
            \gls{sym-Z-n-i1} (\gls{sym-z})$
        &
        $\displaystyle\sum_{\gls{sym-m}}\left(\gls{sym-C-1m-i2}\,
            \gls{sym-R-1m-i2} (\gls{sym-r})
            +
            \gls{sym-C-2m-i2}\,
            \gls{sym-R-2m-i2} (\gls{sym-r})
        \right)
        \times~
        \gls{sym-Z-m-i2} (\gls{sym-z})$
        &
        $\displaystyle\sum_{\gls{sym-k-index}}
            \gls{sym-B-k-e}
            \gls{sym-Lambda-k-e} (\gls{sym-r})
            \gls{sym-Z-k-e} (\gls{sym-z})$ \\
    Partic. potential $\gls{sym-phi-p}(\gls{sym-r},\gls{sym-z})$
        &
        $\displaystyle
        \frac{1}{2 (\gls{sym-h}-\gls{sym-d-1})}
        \left[
            (\gls{sym-z}+\gls{sym-h})^2
            -
            \frac{\gls{sym-r}^2}{2}
        \right]$
        &
        $\displaystyle
        \frac{1}{2 (\gls{sym-h}-\gls{sym-d-2})}
        \left[
            (\gls{sym-z}+\gls{sym-h})^2
            -
            \frac{\gls{sym-r}^2}{2}
        \right]$
        &
        $0$
        \\
        \hline 
    Radial eigen-function $\gls{sym-R-r}$
        &  $\gls{sym-R-1n-i1} (\gls{sym-r})
            =
            \begin{cases}
                \frac12
                    & \gls{sym-n}=0\\[1em]
                \displaystyle
                \frac{
                    \gls{sym-mathrm-I-0}
                    (\gls{sym-lambda-n-i1}\gls{sym-r})
                }{
                    \gls{sym-mathrm-I-0}
                    (\gls{sym-lambda-n-i1}\gls{sym-a-2})
                }
                    & \gls{sym-n}\ge1
            \end{cases}$
        &  \shortstack{
            $
            \gls{sym-R-1m-i2} (\gls{sym-r})
            =
            \begin{cases}
                \frac12
                    & \gls{sym-m}=0\\[1em]
                \displaystyle
                \frac{
                    \gls{sym-mathrm-I-0}
                    (\gls{sym-lambda-m-i2}\gls{sym-r})
                }{
                    \gls{sym-mathrm-I-0}
                    (\gls{sym-lambda-m-i2}\gls{sym-a}_{2})
                }
                    & \gls{sym-m}\ge1
            \end{cases}
            $
            \\
            $
            \gls{sym-R-2m-i2} (\gls{sym-r})
            =
            \begin{cases}
                \frac12
                \ln\!\left(\frac{\gls{sym-r}}
                        {\gls{sym-a-2}}
                \right)
                    & \gls{sym-m}=0\\
                \displaystyle
                \frac{
                    \gls{sym-mathrm-K-0}
                    (\gls{sym-lambda-m-i2}\gls{sym-r})
                }{
                    \gls{sym-mathrm-K-0}
                    (\gls{sym-lambda-m-i2}\gls{sym-a}_{2})
                }
                    & \gls{sym-m}\ge1
            \end{cases}
            $
            }
        &
        $\gls{sym-Lambda-k-e} (\gls{sym-r})
            =
            \begin{cases}
                \displaystyle
                \frac{
                    \gls{sym-mathrm-H-0-1}
                    (\gls{sym-m-0}\gls{sym-r})
                }{
                    \gls{sym-mathrm-H-0-1}
                    (\gls{sym-m-0}\gls{sym-a-2})
                }
                    & \gls{sym-k-index}=0\\[1em]
                \displaystyle
                \frac{
                    \gls{sym-mathrm-K-0}
                    (\gls{sym-m-k}\gls{sym-r})
                }{
                    \gls{sym-mathrm-K-0}
                    (\gls{sym-m-k}\gls{sym-a-2})
                }
                    & \gls{sym-k-index}\ge1
            \end{cases}$
            \\
            \hline 
    Vertical eigen-function $\gls{sym-Z}(\gls{sym-z})$
        &
        $\gls{sym-Z-n-i1} (\gls{sym-z})
        =
        \sqrt{2}\times
        \begin{cases}
            1/\sqrt{2}
                & \gls{sym-n}=0\\[1em]
            \cos\!\left(\gls{sym-lambda-n-i1}
                (\gls{sym-z}+\gls{sym-h})
            \right)
                & \gls{sym-n}\ge1
        \end{cases}$
        &
        $\gls{sym-Z-m-i2} (\gls{sym-z})
        =
        \sqrt{2}\times
        \begin{cases}
            1/\sqrt{2}
                & \gls{sym-m}=0\\[1em]
            \cos\!\left(\gls{sym-lambda-m-i2}
                (\gls{sym-z}+\gls{sym-h})
            \right)
                & \gls{sym-m}\ge1
        \end{cases}$
        &
        $\gls{sym-Z-k-e} (\gls{sym-z})
        =
        \gls{sym-N-k}^{-\frac12}
        \times
        \begin{cases}
            \cosh\!\left(\gls{sym-m-0}
                (\gls{sym-z}+\gls{sym-h})
            \right)
                & \gls{sym-k-index}=0\\[1em]
            \cos\!\left(\gls{sym-m-k}
                (\gls{sym-z}+\gls{sym-h})
            \right)
                & \gls{sym-k-index}\ge1
        \end{cases}$
        \\
        \hline
    Eigen-value& 
        $\displaystyle
        \gls{sym-lambda-n-i1}
        =
        \frac{\gls{sym-n}\gls{sym-pi}}
            {\gls{sym-h}-\gls{sym-d-1}},
        \quad
        \gls{sym-n}\ge1$
        &
        $\displaystyle
        \gls{sym-lambda-m-i2}
        =
        \frac{\gls{sym-m}\gls{sym-pi}}
            {\gls{sym-h}-\gls{sym-d-2}},
        \quad
        \gls{sym-m}\ge1$
        &
        \ifdefined\DISSERTATION
            $\gls{sym-m-k}$ (see \Cref{aor:eq:dispersion})
        \else
            $\displaystyle
            \dispersion
            $
        \fi
    \\
    \hline
    \endtabular
    \caption{Equations for potential (homogeneous and particular) and eigenfunctions (radial and vertical) for each region.}
    \label{tab:MEEM-eigenfunctions}
\end{table}
\ifdefined\DISSERTATION
    The dispersion relation is a nonlinear equation that implicitly gives the exterior-region eigenvalues $\gls{sym-m-k}$:
    \begin{equation}
        \label{aor:eq:dispersion}
    \dispersion
    \end{equation}
\fi
In \Cref{tab:MEEM-eigenfunctions},  $\gls{sym-mathrm-I-0}$,  $\gls{sym-mathrm-K-0}$, and $\gls{sym-mathrm-H-0-1}$ are different Bessel functions of order zero; 1M and 2M mean body 1 and 2 (spar and float) are moving respectively, while 1S and 2S mean each is stationary; and the $\gls{sym-N-k}$ expression is defined as:

\begin{equation}
    \gls{sym-N-k} = \frac{1}{2}\left(1+\frac{\gls{sym-f-k}}{2\gls{sym-m-k}\gls{sym-h}}  \right)~
    \textrm{ where }
    \gls{sym-f-k} =
    \begin{cases}
        \sinh(2\gls{sym-m-0}\gls{sym-h}), & \gls{sym-k-index}=0 \\ \sin(2\gls{sym-m-k}\gls{sym-h}), & \gls{sym-k-index}\geq1
    \end{cases}
\end{equation}

 %%%%%%%%%%%%%%%%%%%%%%%%%%%%%%%%%%%%%%%%%%%%%%%%%%%%%%%%%%%%%%% begin equation helper functions %%%%%%%%%%%%%%%%%%%%%%%%%%%%%%%%%%%%%%%%%%%%%%%%%%%%%%%%%%%
%% integral definitions
\newcommand{\RintOneDefn}{
    \shortstack{
        $\displaystyle
            \gls{sym-boldsymbol-mathcal-R}_{1j} =
            \int\limits_{\gls{sym-a-in}}^{\gls{sym-a-out}} \gls{sym-vec-R}_{1j} (\gls{sym-r})\gls{sym-r} dr$
        \\
            for $\gls{sym-j}= (\gls{sym-n},\gls{sym-m})
        $
    }
}

\newcommand{\RintTwoDefn}{
    $\displaystyle \gls{sym-boldsymbol-mathcal-R}_{2\gls{sym-m}} = \int\limits_{\gls{sym-a-in}}^{\gls{sym-a-out}} \gls{sym-vec-R}_{2\gls{sym-m}}(\gls{sym-r})rdr$
}

\newcommand{\ZmnDefn}{
    \shortstack{
        $\displaystyle \gls{sym-boldsymbol-mathcal-Z}_{nm} = \gls{sym-boldsymbol-mathcal-Z}_{mn}^\top =$
        \\
        $\displaystyle \int_{-\gls{sym-h}}^{-\gls{sym-d-1}} \gls{sym-vec-Z-n-i1}^{\top}\gls{sym-vec-Z-m-i2} dz$
    }
}
\newcommand{\ZmkDefn}{
    \shortstack{
        $\displaystyle \gls{sym-boldsymbol-mathcal-Z}_{mk} = \gls{sym-boldsymbol-mathcal-Z}_{km}^\top = $
        \\
        $\displaystyle \int_{-\gls{sym-h}}^{-\gls{sym-d-2}} \gls{sym-vec-Z-m-i2}^{\top} \gls{sym-vec-Z-k-e} dz $
    }
}

%% R definitions
\newcommand{\RintOneJzero}{
    $\displaystyle\frac{\gls{sym-a-out}^2-\gls{sym-a-in}^2}{4}$
}

\newcommand{\RintOneJOne}{
    $\displaystyle
        \frac
            {\gls{sym-a-out} \gls{sym-mathrm-I-1} (\gls{sym-vec-lambda}_{\gls{sym-j}} \gls{sym-a-out}) - \gls{sym-a-in}\gls{sym-mathrm-I-1} (\gls{sym-vec-lambda}_{\gls{sym-j}} \gls{sym-a-in})}
            {\gls{sym-vec-lambda}_{\gls{sym-j}} \gls{sym-mathrm-I-0} (\gls{sym-vec-lambda}_{\gls{sym-j}} \gls{sym-a-2})}
    $
}

\newcommand{\RintTwoJzero}{
    $\displaystyle
        \frac
        {
            \gls{sym-f}_{2m}\gls{sym-a-in}^2-\gls{sym-a-out}^2
        }
        {8}
    $
    where $
        \gls{sym-f}_{2m} = 1 + 2\ln\frac{\gls{sym-a-out}}{\gls{sym-a-in}}
    $
}

\newcommand{\RintTwoJOne}{
    $\displaystyle
    \frac
    {
        \gls{sym-a-in}\,\gls{sym-mathrm-K-1} (\gls{sym-vec-lambda-m} \gls{sym-a-in})
        -\gls{sym-a-out}\,\gls{sym-mathrm-K-1} (\gls{sym-vec-lambda-m} \gls{sym-a-out})
    }
    {
        \gls{sym-vec-lambda-m} \,\gls{sym-mathrm-K-0} (\gls{sym-vec-lambda-m} \gls{sym-a-2})
    }
    $
}

% Znm formulas
\newcommand{\ZnZeroMZero}{
    $\gls{sym-h}-\gls{sym-d-1}$
}

\newcommand{\ZnZeroMOne}{
    $\displaystyle
     \frac
                {\sqrt{2}\sin\left(
				\gls{sym-vec-lambda-m} (\gls{sym-h} - \gls{sym-d-1})
			\right)}
                { \gls{sym-vec-lambda-m}}
    $
}

\newcommand{\ZnOneMZero}{
    $\displaystyle 0
     % this more complicated formula is only needed for multi-region
     %\frac
      %          {\sqrt{2}\sin (\gls{sym-vec-lambda-n}^{\top} (\gls{sym-h} - \gls{sym-d-1}))}
        %        { \gls{sym-vec-lambda-n}^{\top}}
    $
}

\newcommand{\ZnOneMOne}{
$\displaystyle
\begin{cases}
\displaystyle\frac{
            2 (-1)^{\gls{sym-vec-n}}
            }
            {
                    1 - \left(\frac{\gls{sym-vec-lambda-n}^{\top}}{  \gls{sym-vec-lambda-m}}
                        \right)^2
            }
\displaystyle\frac{
            \sin\left(
                \gls{sym-vec-lambda-m} (\gls{sym-h}-\gls{sym-d-1})
         	\right)
            }
            {
            \gls{sym-vec-lambda-m}
            }
,&  \gls{sym-lambda-n-i1} \neq \gls{sym-lambda-m-i2}\\
\gls{sym-h} -\gls{sym-d-1}, & \gls{sym-lambda-n-i1} = \gls{sym-lambda-m-i2}
\end{cases}
$
% this more complicated formula is only needed for multi-region
% \shortstack{
%     $\displaystyle
%     \frac
%         {\sin\left(% 		(\gls{sym-vec-lambda-n}^{\top} + \gls{sym-vec-lambda-m}) (\gls{sym-h}-\gls{sym-d-1})
% 	\right)}
%         {\gls{sym-vec-lambda-n}^{\top} + \gls{sym-vec-lambda-m}}
%     + \gls{sym-Z}_{nm}^-$
%     \\
%     where $\gls{sym-Z}_{nm}^-=
%     \begin{cases}
%         \frac
%             {\sin\left(% 		(\gls{sym-vec-lambda-n}^{\top} - \gls{sym-vec-lambda-m}) (\gls{sym-h}-\gls{sym-d-1})
% 	\right)}
%             {(\gls{sym-vec-lambda-n}^{\top} - \gls{sym-vec-lambda-m})},
%         & \gls{sym-lambda-n-i1} \neq \gls{sym-lambda-m-i2}
%         \\
%         \gls{sym-h}-\gls{sym-d-1},
%         & \gls{sym-lambda-n-i1} = \gls{sym-lambda-m-i2}
%     \end{cases}
%   $
%  }
}

% Zmk formulas
\newcommand{\ZmZeroKZero}{
    $\displaystyle \frac{\sinh(\gls{sym-m-0}(\gls{sym-h}-\gls{sym-d-2}))}{\gls{sym-m-0} \sqrt{\gls{sym-N-0}}}$
}

\newcommand{\ZmOneKZero}{
    $\displaystyle \frac{\sqrt{2}~\gls{sym-m-0} (-1)^{\vec{\gls{sym-m}}} \sinh(\gls{sym-m-0}(\gls{sym-h}-\gls{sym-d-2}))}{\sqrt{\gls{sym-N-0}} \left(\gls{sym-m-0}^2+{(\gls{sym-vec-lambda-m}^{\top})}^2\right)}$
}

\newcommand{\ZkOneMZero}{
    $\displaystyle\frac{
                        \sin (\gls{sym-vec-m-k} (\gls{sym-h}-\gls{sym-d-2}))
                        }
                        {
                        \gls{sym-vec-m-k}
                        \sqrt{\vec{\gls{sym-N-k}}}
                        }$
}

\newcommand{\ZkOneMOne}{
    $\displaystyle
    \begin{cases}
            \displaystyle\frac{1}{\sqrt{2\vec{\gls{sym-N-k}}}}
            \left(
                \displaystyle\frac
                    {
                        \sin\left((\gls{sym-h}-\gls{sym-d-2})
                            (\gls{sym-vec-m-k} + \gls{sym-vec-lambda-m}^{\top})
                        \right)
                    }
                    {
                        \gls{sym-vec-m-k}
                        + \gls{sym-vec-lambda-m}^{\top}
                    }
            + \right.
                 \\
                \left.~~~~~~~~~
                \displaystyle\frac
                    {\sin\left((\gls{sym-h}-\gls{sym-d-2}) (\gls{sym-vec-m-k}- \gls{sym-vec-lambda-m}^{\top})\right)}
                    {\gls{sym-vec-m-k}- \gls{sym-vec-lambda-m}^{\top}}
            \right),
        & |\gls{sym-m-k}| \neq \gls{sym-lambda-m}
        \\
        \displaystyle\frac{\gls{sym-h}-\gls{sym-d-2}}{2},
        & |\gls{sym-m-k}| = \gls{sym-lambda-m}
    \end{cases}$
}

 %%%%%%%%%%%%%%%%%%%%%%%%%%%%%%%%%%%%%%%%%%%%%%%%%%%%%%%%%%%%%%% end equation helper functions %%%%%%%%%%%%%%%%%%%%%%%%%%%%%%%%%%%%%%%%%%%%%%%%%%%%%%%%%%%

\Cref{tab:meem-integrals} lists several integrals of the radial and vertical eigenfunctions, $\gls{sym-boldsymbol-mathcal-R}$ and $\boldsymbol{\mathcal{\gls{sym-Z}}}$ respectively, that will be needed in the calculations to follow.

\begin{table}[htbp]
    \centering
    \caption{Eigenfunction integrals}
    \label{tab:meem-integrals}
    \begin{tabular}{|M{0.16\linewidth}|M{0.05\linewidth}|M{0.20\linewidth}|M{0.48\linewidth}|} \hline 
        \multicolumn{2}{|M{0.21\linewidth}|}{Integral}           & $\gls{sym-m},\gls{sym-j}=0$                       & $\gls{sym-m},\gls{sym-j}\geq1$            \\ \hline 
        \multicolumn{2}{|M{0.21\linewidth}|}{\RintOneDefn}       & \RintOneJzero                 & \RintOneJOne          \\ \hline 
        \multicolumn{2}{|M{0.21\linewidth}|}{\RintTwoDefn}       & \RintTwoJzero                 & \RintTwoJOne          \\ \hline 
        \multirow{6}{*}{\ZmnDefn}   & \multirow{2}{*}{$\gls{sym-n}=0$}     & \multirow{2}{*}{\ZnZeroMZero} & \ZnZeroMOne           \\ \cline{2-4} 
                                    & \multirow{3}{*}{$\gls{sym-n}\geq1$}  & \multirow{3}{*}{\ZnOneMZero}  & \ZnOneMOne            \\ \hline
        \multirow{9}{*}{\ZmkDefn}   & \multirow{2}{*}{$\gls{sym-k-generic-index}=0$}     & \multirow{2}{*}{\ZmZeroKZero} & \ZmOneKZero           \\ \cline{2-4}
                                    & \multirow{4}{*}{$\gls{sym-k-generic-index}\geq1$}  & \multirow{4}{*}{\ZkOneMZero}  & \ZkOneMOne            \\ \hline
    \end{tabular}
\end{table}


\subsection{Matching Across Fluid Boundaries}
\label{sec:appendix-meem-matching}

The eigencoefficients must be selected to enforce the radial velocity body boundary condition and the matching of the potentials and radial velocities at the edges of each region, earning this technique the name \gls{acr-meem}.
The radiation problem was first solved this way for a floating cylinder in 1980 \citep{yeung_added_1981}.

First, the infinite sums in $\gls{sym-phi-h}$ must be truncated.
Assuming truncation to $\gls{sym-N}$ terms in \textit{i1}, $\gls{sym-M}$ terms in \textit{i2}, and $\gls{sym-K}$ terms in \textit{e}, the total number of eigencoefficients to solve for is $\gls{sym-N}+2\gls{sym-M}+\gls{sym-K}$.
For a 3-region problem, there are 2 boundaries.
Thus there are four matching equations: (1) potential at $\gls{sym-a-1}$, (2) potential at $\gls{sym-a-2}$, (3) velocity at $\gls{sym-a-1}$, and (4) velocity at $\gls{sym-a-2}$.
As-is, this is not enough equations ($4 < \gls{sym-N}+2\gls{sym-M}+\gls{sym-K}$).
We must leverage eigenfunction orthogonality to get enough equations.
The first equation will turn into $\gls{sym-N}$ equations; the second and third each give $\gls{sym-M}$; the fourth $\gls{sym-K}$.
The transformation uses the following property of orthogonality.
Consider a generic function $\gls{sym-Y}(\gls{sym-x})$ expressed as a series with coefficients $\gls{sym-alpha-generic}$ and basis functions $e(\gls{sym-x})$: $\gls{sym-Y}(\gls{sym-x})=\sum_i \gls{sym-alpha-generic}_i \gls{sym-e-i}(\gls{sym-x})$.
If $\gls{sym-e-j}(\gls{sym-x})$ is orthogonal to $\gls{sym-e-i}(\gls{sym-x})$ from $\gls{sym-x} = \gls{sym-a}$ to $\gls{sym-b}$, then:

\begin{equation}
\begin{aligned}
       & \int\limits_{\gls{sym-a-b}} \gls{sym-Y} (\gls{sym-x})\gls{sym-e-j} (\gls{sym-x})dx = (\gls{sym-b}-\gls{sym-a}) <\gls{sym-Y},\gls{sym-e-j}>= (\gls{sym-b}-\gls{sym-a})<\sum_i \gls{sym-alpha-generic}_i \gls{sym-e-i}, \gls{sym-e-j}> \\ &= (\gls{sym-b}-\gls{sym-a}) \sum_i \gls{sym-alpha-generic}_i <\gls{sym-e-i},\gls{sym-e-j}> = (\gls{sym-b}-\gls{sym-a}) \sum_i \gls{sym-alpha-generic}_i \gls{sym-delta-ij} = (\gls{sym-b}-\gls{sym-a}) \gls{sym-alpha-generic}_{\gls{sym-j}}
\end{aligned}
\end{equation}

where $<\cdot,\cdot>$ is the inner product and $\gls{sym-delta-ij}$ is Kronecker's delta.
In the current hydrodynamics problem, the basis functions are the vertical eigenfunctions $\gls{sym-Z-n-i1}$, $\gls{sym-Z-m-i2}$, and $\gls{sym-Z-k-e}$.
Orthogonality of each eigenfunction can be verified with the inner product.
In the first region, for example, $<\gls{sym-Z-n-1-i-1},Z_{n_2}^{i_1}>=\delta_{n_1 n_2}$.
Note that eigenfunctions of different domains are not orthogonal, and their inner products will be expressed as coupling integrals in \Cref{tab:meem-integrals}.

For each of the four matching equations, the property of orthogonality applies only after multiplying by the appropriate eigenfunction and integrating over appropriate bounds.
For the potential matching equations, multiply both sides by the eigenfunction of the region with smaller fluid height (so $\gls{sym-Z-n-i1}$ at $\gls{sym-a-1}$ and $\gls{sym-Z-m-i2}$ at $\gls{sym-a-2}$).
Then integrate over that fluid height ($\gls{sym-z}=-\gls{sym-h}$ to $-\gls{sym-d-1}$ at $\gls{sym-a-1}$, and $-\gls{sym-h}$ to $-\gls{sym-d-2}$ at $\gls{sym-a-2}$).
For velocity matching, multiply instead by the eigenfunction corresponding to the larger region, while still integrating over the smaller region.
In velocity matching, an extra step is required to incorporate the boundary condition of zero radial velocity along the radial surface of the body.
Since it is zero-valued, the integral of this velocity may be added to one side of the equation (the one corresponding to the velocity of the larger region) to change the integration bounds only on that side.
This manifests in the bounds of the coupling integrals to be presented in \Cref{tab:meem-integrals}.
Other combinations of eigenfunction multiplication or integration besides those described above are not useful since they result in integrating a quantity on a region where it is undefined, or a form unsuitable for the application of the orthogonality property.

\subsection{Block Matrix Structure}

Once orthogonality is applied, the matching equations create a linear system $\gls{sym-mathbf-A}\gls{sym-vec-x}=\gls{sym-vec-b}$ where $\gls{sym-mathbf-A}$ is a complex sparse $(\gls{sym-N}+2\gls{sym-M}+\gls{sym-K})$x$(\gls{sym-N}+2\gls{sym-M}+\gls{sym-K})$ square matrix corresponding to the homogeneous case, $\gls{sym-vec-x}=[\gls{sym-vec-C-1n-i1}, \gls{sym-vec-C-1m-i2}, \gls{sym-vec-C-2m-i2}, \vec{\gls{sym-B}_{\gls{sym-k-index}}^{e}}]$ is the complex eigencoefficient vector, and $\gls{sym-vec-b}$ is the real boundary condition vector corresponding to the particular case.
We elaborate on the block structure of the A-matrix and b-vector, an implementation detail that prior discussion of \gls{acr-meem} overlooks.
The A-matrix and b-vector block structures are shown in \Cref{tab:MEEM-A-matrix,tab:MEEM-b-vector} respectively.
They are written in compact notation using row vectors of basis functions, so $\gls{sym-vec-R-1n-i1}=[R_{10}^{i1}, R_{11}^{i1},\dots, R_{1(N-1)}^{i1}]$ and so on.
Each basis function is evaluated at the radius described to the left of its row in the table. $0_{ij}$ and $1_{ij}$ are the $i$ x $\gls{sym-j}$ matrices of zeros and ones respectively; diag($\cdot$) constructs a diagonal matrix from a vector; and $\odot$ is the Hadamard (element-wise) product.

\begin{table}[htbp]
    \centering
    \caption{\gls{acr-meem} A-matrix}
    \label{tab:MEEM-A-matrix}
    \begin{tabular}{|M{0.12\linewidth}|c||c|c|c|c|}\hline
 & & $\gls{sym-vec-C-1n-i1}$& $\gls{sym-vec-C-1m-i2}$& $\gls{sym-vec-C-2m-i2}$&$\gls{sym-vec-B-k-e}$\\\hline
          &size&  N&  M&  M& K
          \\ \hline \hline
            \shortstack{
                            $\gls{sym-phi-i1}=\gls{sym-phi-i2}$
                        \\
                            at $\gls{sym-r}=\gls{sym-a-1}$
                        }
            &N
            &  $(\gls{sym-h}-\gls{sym-d-1})~\gls{sym-mathrm-diag}(\gls{sym-vec-R-1n-i1})$
            &  $-\gls{sym-boldsymbol-mathcal-Z}_{nm}\odot 1_{\gls{sym-N}1}\gls{sym-vec-R-1m-i2}$
            &  $-\gls{sym-boldsymbol-mathcal-Z}_{nm}\odot 1_{\gls{sym-N}1}\gls{sym-vec-R-2m-i2}$& $0_{NK}$
          \\ \hline
            \shortstack{
                            $\gls{sym-phi-i2}=\gls{sym-phi-e}$
                        \\
                            at $\gls{sym-r}=\gls{sym-a-2}$
                        }
            &M
            &  $0_{MN}$&  $(\gls{sym-h}-\gls{sym-d-2})~\gls{sym-mathrm-diag}(\gls{sym-vec-R-1m-i2})$
            &  $(\gls{sym-h}-\gls{sym-d-2})~\gls{sym-mathrm-diag}(\gls{sym-vec-R-2m-i2})$
            & $-\gls{sym-boldsymbol-mathcal-Z}_{mk}\odot 1_{\gls{sym-M}1}\vec{\Lambda}_{\gls{sym-k-index}}$
          \\ \hline
            \shortstack{
                            $\frac{\partial}{\partial \gls{sym-r}}
                            \gls{sym-phi-i1}=\frac{\partial}{\partial \gls{sym-r}}
                            \gls{sym-phi-i2}$
                        \\
                            at $\gls{sym-r}=\gls{sym-a-1}$
                        }
            &M
            &  $- \gls{sym-boldsymbol-mathcal-Z}_{mn} \odot 1_{\gls{sym-M}1} \frac{\partial}{\partial \gls{sym-r}}\gls{sym-vec-R-1n-i1}$
            &  $(\gls{sym-h}-\gls{sym-d-2})~\gls{sym-mathrm-diag}(\frac{\partial}{\partial \gls{sym-r}}\gls{sym-vec-R-1m-i2})$
            &  $(\gls{sym-h}-\gls{sym-d-2})~\gls{sym-mathrm-diag}(\frac{\partial}{\partial \gls{sym-r}}\gls{sym-vec-R-2m-i2})$
            & $0_{MK}$
          \\ \hline
            \shortstack{
                            $\frac{\partial}{\partial \gls{sym-r}}
                            \gls{sym-phi-i2}=\frac{\partial}{\partial \gls{sym-r}}
                            \gls{sym-phi-e}$
                        \\
                            at $\gls{sym-r}=\gls{sym-a-2}$
                        }
                &K
            &  $0_{KN}$
            &  $-\gls{sym-boldsymbol-mathcal-Z}_{km} \odot 1_{\gls{sym-K}1}\frac{\partial}{\partial \gls{sym-r}}\gls{sym-vec-R-1m-i2}$
            &  $-\gls{sym-boldsymbol-mathcal-Z}_{km} \odot 1_{\gls{sym-K}1}\frac{\partial}{\partial \gls{sym-r}}\gls{sym-vec-R-2m-i2}$
            & $\gls{sym-h}~\gls{sym-mathrm-diag}
            (\frac{\partial}{\partial \gls{sym-r}} \gls{sym-vec-Lambda-k-e})$
        \\ \hline
    \end{tabular}
\end{table}

\begin{table}[htbp]
    \centering
    \caption{\gls{acr-meem} b-vector}
    \label{tab:MEEM-b-vector}
    \begin{tabular}{|c|c|} \hline
         N& $\displaystyle
\int_{-\gls{sym-h}}^{-\gls{sym-d-1}} (\gls{sym-phi-p}^{i2} - \gls{sym-phi-p}^{i1})\gls{sym-vec-Z-n-i1}^{~\top} dz
$ \\ \hline
         M& $\displaystyle-
\int_{-\gls{sym-h}}^{-\gls{sym-d-2}} \gls{sym-phi-p}^{i2} \gls{sym-vec-Z-m-i2}^{~\top} dz
$\\ \hline
         M& $\displaystyle
\int_{-\gls{sym-h}}^{-\gls{sym-d-1}} \frac{\partial}{\partial \gls{sym-r}}\gls{sym-phi-p}^{i1}\gls{sym-vec-Z-m-i2}^{i2~\top} dz - \int_{-\gls{sym-h}}^{-\gls{sym-d-2}}\frac{\partial}{\partial \gls{sym-r}}\gls{sym-phi-p}^{i2}\gls{sym-vec-Z-m-i2}^{~\top} dz
$\\ \hline
         K& $\displaystyle
\int_{-\gls{sym-h}}^{-\gls{sym-d-2}} \frac{\partial}{\partial \gls{sym-r}}\gls{sym-phi-p}^{i2} \vec{\gls{sym-Z}}_{\gls{sym-k-index}}^{e~\top} dz
$\\ \hline
    \end{tabular}
\end{table}

The dense blocks contain coupling integrals $\gls{sym-boldsymbol-mathcal-Z}$ of the vertical eigenfunctions.

Of the sixteen blocks that make up the matrix, six are diagonal, four are zero, and six are dense, resulting in the sparsity pattern shown in \Cref{fig:sparsity}.
An even sparser matrix could be obtained with the alternate eigenfunction scaling for the second region described in \citet{chau_inertia_2012}.
\begin{figure}[htbp]
    \centering
    \includegraphics[\FigureOptions{aor:figs/from-matlab/meem_sparsity.pdf}]{figs/from-matlab/meem_sparsity.pdf}
    \caption{A-matrix sparsity pattern, shown for $\gls{sym-N}=\gls{sym-M}=\gls{sym-K}=4$}
    \label{fig:sparsity}
\end{figure}


\subsection{Calculation of Outputs}

Once the eigencoefficients have been calculated by solving the linear system for $\gls{sym-vec-x}$, the hydrodynamic radiation coefficients $\gls{sym-A-h}$ and $\gls{sym-B-h}$ are found by integrating the potentials and velocities over the body surface:
\begin{alignat}{5}\label{eq:meem-c-derivation}
    \gls{sym-A-h} + \frac{i\gls{sym-B-h}}{\gls{sym-omega}}&=\gls{sym-rho-w} \gls{sym-h}^3  \int_{\gls{sym-theta}=0}^{2\gls{sym-pi}}
    & \int_{\gls{sym-r}=\gls{sym-a-in}}^{\gls{sym-a-out}}
    &\gls{sym-phi} (\gls{sym-r},\gls{sym-z}) \frac{\partial \gls{sym-phi} (\gls{sym-r},\gls{sym-z})}{\partial \gls{sym-z}}
    \gls{sym-r} d\gls{sym-r} d\gls{sym-theta}  \nonumber \\
    %
    %
    &= 2\pi \rho h^3
    & \int_{r=a_{in}}^{a_{out}}&
   \left(\phi_p (r,z) +
         \sum_{j_1} C_{j_1} R_{j_1} (r)Z_{j_1} (z)
   \right) \nonumber \\
   %
  & & &  \left(\frac{\partial \phi_p (r,z)}{\partial z}
               + \sum_{j_2} C_{j_2} R_{j_2} (r)
               \frac{\partial Z_{j_2} (z)} {\partial z}
        \right)r dr \nonumber \\
    %
    %
    &= 2\pi \rho h^3 & \left[
    % first term
    \int_{r=a_{in}}^{a_{out}} \right. &
    \left. \phi_p (r,z)
        \frac{\partial \phi_p (r,z)}{\partial z} rdr
    \right.
    \\
    % second term
    & &  + \sum_{j_1} &C_{j_1} Z_{j_1} (z)
    \int_{r=a_{in}}^{a_{out}}R_{j_1} (r)
    \frac{\partial \phi_p (r,z)}{\partial z}rdr \nonumber
    \\
    % third term
    &&   + \sum_{j_2}&
    \frac{\partial Z_{j_2} (z)}{\partial z} C_{j_2}
    \int_{r=a_{in}}^{a_{out}}
    \phi_p (r,z)  R_{j_2} (r)rdr \nonumber
    \\
    % fourth term
    & & +  \sum_{j_1}&
    \left.
        C_{j_1}Z_{j_1} (z) \sum_{j_2} C_{j_2}
        \frac{\partial Z_{j_2} (z)}{\partial z}
        \int_{r=a_{in}}^{a_{out}}
        R_{j_1} (r)R_{j_2} (r)rdr
    \right] \nonumber
\end{alignat}
where the generic indices $\gls{sym-j}_1$ and $\gls{sym-j}_2$ are used to represent $\gls{sym-n}$ or $\gls{sym-m}$ depending on the region.
For regions with multiple radial eigenfunctions or surfaces covering multiple fluid regions, summation over each eigenfunction and region respectively is implied.
After substituting the eigenfunctions from \Cref{tab:MEEM-eigenfunctions} into \Cref{eq:meem-c-derivation}, all $\gls{sym-z}$-dependent quantities are evaluated at the body draft on the integration surface, $\gls{sym-z}=-d$.
This means $\frac{\partial \gls{sym-phi-p}}{\partial \gls{sym-z}}=1$, and $\gls{sym-Z-j}(\gls{sym-z})$ simplifies to
\begin{equation}
    \gls{sym-Z-j} (\gls{sym-z}=-d) =
    \begin{cases}
        1,& \gls{sym-j}=0
        \\
        \sqrt2(-1)^{\gls{sym-j}},& \gls{sym-j} \geq 1
    \end{cases}
\end{equation}
The third and fourth integrals of \Cref{eq:meem-c-derivation} vanish, and the second integral is the radial eigenfunction integral $\gls{sym-boldsymbol-mathcal-R}$ expressed previously in \Cref{tab:meem-integrals}.

The first and second term are independent of and scale linearly with the eigencoefficients $\gls{sym-vec-x}$, respectively.
The radiation coefficients are thus computed as
\begin{equation}\label{eq:hydro}
    \gls{sym-A-h} + \frac{i\gls{sym-B-h}}{\gls{sym-omega}}=2\gls{sym-pi} \gls{sym-rho-w} \gls{sym-h}^3 (\gls{sym-c-0} + \gls{sym-vec-c}\cdot\gls{sym-vec-x}) =2\gls{sym-pi} \gls{sym-rho-w} \gls{sym-h}^3 (\gls{sym-c-0} + \gls{sym-vec-c} \cdot \gls{sym-mathbf-A}^{-1} \gls{sym-vec-b})
\end{equation}
with output constant $\gls{sym-c-0}$ and output row vector $\gls{sym-vec-c}$ defined in \Cref{eq:c0,tab:MEEM-c-vector} respectively:
\begin{equation}\label{eq:c0}
    \gls{sym-c-0} =  \frac{\left({\gls{sym-a-out}}-{\gls{sym-a-in}}\right)\,\left(-{\gls{sym-a-in}}-{\gls{sym-a-out}}+4 (\gls{sym-h}-d)^2\right)}{16\,\left(\gls{sym-h} - d\right)}
\end{equation}
\begin{table}[h]
    \centering
    \caption{Output vector $\gls{sym-vec-c}$}
    \label{tab:MEEM-c-vector}
\begin{tabular}{|l|c|c|c|c|} \hline
          Body&N&  M&  M& K\\ \hline
          Float&$0_{1\gls{sym-N}}$&  $\gls{sym-vec-Z-m-i2}(\gls{sym-z}=-d) \odot\gls{sym-boldsymbol-mathcal-R}_{1\gls{sym-m}}$&  $\gls{sym-vec-Z-m-i2}(\gls{sym-z}=-d) \odot\gls{sym-boldsymbol-mathcal-R}_{2\gls{sym-m}}$& $0_{1\gls{sym-K}}$\\ \hline
          Spar& $\gls{sym-vec-Z-n-i1}(\gls{sym-z}=-d) \odot\gls{sym-boldsymbol-mathcal-R}_{1\gls{sym-n}}$& $0_{1\gls{sym-M}}$& $0_{1\gls{sym-M}}$&$0_{1\gls{sym-K}}$\\\hline
    \end{tabular}
    \end{table}
The appropriate dimensions are
\begin{equation}
    \gls{sym-a-out} = \begin{cases}
        \gls{sym-a-2}, & \text{float} \\
        \gls{sym-a-1}, & \text{spar}
    \end{cases}
    \qquad
    \gls{sym-a-in} = \begin{cases}
        \gls{sym-a-1}, & \text{float} \\
        0, & \text{spar}
    \end{cases}
    \qquad
   d = \begin{cases}
        -\gls{sym-d-2}, & \text{float} \\
        -\gls{sym-d-1}, & \text{spar}
    \end{cases}
\end{equation}
to obtain the hydrodynamic coefficients of the float or the spar respectively.

Using the A-matrix and b- and c-vectors, the \gls{acr-meem} solution directly calculates radiation coefficients $\gls{sym-B-h}$ and $\gls{sym-A-h}$.
The hydrostatic stiffness $\gls{sym-K-h}$ and the excitation coefficient $\gls{sym-gamma}$ are calculated as follows:
\begin{equation}\label{eq:gamma-K}
    |\gls{sym-gamma}|  = \sqrt{\frac{ 4 \gls{sym-rho-w} \gls{sym-g} \gls{sym-V-g}  \gls{sym-B-h}} {\gls{sym-m-0}}}, \quad % excitation
   \angle \gls{sym-gamma} = -\frac{\gls{sym-pi}}{2} + \angle\frac{ {\gls{sym-B-k-e}}_{=0}}{\gls{sym-mathrm-H-0-1} (\gls{sym-m-0} \gls{sym-a-2})},\quad
    \gls{sym-K-h}       = \gls{sym-rho-w} \gls{sym-g} \underbrace{\frac{\gls{sym-pi}}{4} (\gls{sym-D-f}^2 - \gls{sym-D-s}^2)}_{\gls{sym-A-w-f}},\quad  % hydrostatic stiffness
    \gls{sym-K-s} = \gls{sym-rho-w} \gls{sym-g} \underbrace{\frac{\gls{sym-pi}}{4} \gls{sym-D-s}^2}_{\gls{sym-A-w-s}}
\end{equation}
where $\gls{sym-g}$ is the acceleration due to gravity, $\gls{sym-rho-w}$ is the density of water, $\gls{sym-m-0}$ is the wavenumber, $\gls{sym-V-g}$ is the finite depth group velocity, $\gls{sym-mathrm-H-0-1}$ is the zeroth-order Hankel function of the first kind, and $\gls{sym-A-w}$ is the waterplane area \citep{newman}.
This method of calculating excitation from damping is the well-known Haskind relation.
Note that while the excitation magnitude $|\gls{sym-gamma}|$ depends on the radiation damping $\gls{sym-B-h}$ which in turn depends on all the inner region eigencoefficients ($\gls{sym-vec-C-1m-i2}$ and $\gls{sym-vec-C-2m-i2}$ for float excitation and $\gls{sym-vec-C-1n-i1}$ for the spar excitation), the excitation phase $\angle\gls{sym-gamma}$ depends only on the first exterior eigencoefficient, ${\gls{sym-B-k-e}}_{=0}$.

\subsection{Verification and Validation}

Hydro coefficient results are verified by comparing to a benchmark shallow-water concentric-cylinder \gls{acr-meem} solution in \citet{chau_inertia_2012}.
Excellent agreement is observed, shown in \Cref{fig:meem-yeung-validation}.
The results are also experimentally validated in the studies \citet{chau_inertia_2012,son_performance_2016}.
Previously in \Cref{sec:dynamic-validation}, the N=M=K=11 case was compared to \gls{acr-wamit} results for \gls{acr-rm3} in deep water.

\begin{figure}[htbp]
    \centering
    \includegraphics[\FigureOptions{aor:figs/from-matlab/meem_validation.pdf}]{figs/from-matlab/meem_validation.pdf}
    \caption{Nondimensional added mass and damping coefficient validation against \citep{chau_inertia_2012}}
    \label{fig:meem-yeung-validation}
\end{figure}

\subsection{Convergence}

As $\gls{sym-N},\gls{sym-M},\gls{sym-K}\rightarrow\infty$, matching quality improves, and hydro coefficients converge toward their true values.
Previous \gls{acr-meem} papers use $\gls{sym-N}=\gls{sym-M}=\gls{sym-K}=50$ to obtain 4-digit matching accuracy without elaborating on convergence properties \citep{chau_inertia_2012}.
We observe that potential matching converges faster than velocity matching.
\Cref{fig:meem-matching} shows the matching behavior for $\gls{sym-N}=\gls{sym-M}=\gls{sym-K}=11$, where potential matches well but velocity still has noticeable mismatch.
\begin{figure}[htbp]
    \centering
    \includegraphics[\FigureOptions{aor:figs/from-matlab/meem_matching.pdf}]{figs/from-matlab/meem_matching.pdf}
    \caption{Matching for $\gls{sym-N}=\gls{sym-M}=\gls{sym-K}=11$ for benchmark geometry}
    \label{fig:meem-matching}
\end{figure}

Hydro coefficient convergence depends on the geometry: the benchmark shallow-water geometry of \citep{chau_inertia_2012} converges to within 0.25\% with only $\gls{sym-N}=\gls{sym-M}=\gls{sym-K}=4$, but \gls{acr-rm3} requires $\gls{sym-N}=\gls{sym-M}=\gls{sym-K}>10$, shown in \Cref{fig:meem-convergence}.
There, damping converges well at low frequencies but requires more harmonics at higher frequencies, while added mass has similar convergence across frequencies.
\begin{figure}[htbp]
    \centering
    \includegraphics[\FigureOptions{aor:figs/from-matlab/meem_convergence_vs_omega.pdf}]{figs/from-matlab/meem_convergence_vs_omega.pdf}
    \caption{Convergence for $\gls{sym-N}=\gls{sym-M}=\gls{sym-K}=(5,10,20,30)$ for \gls{acr-rm3}}
    \label{fig:meem-convergence}
\end{figure}


\subsection{Numerical Notes}\label{sec:appendix-meem-numerics}
Note that radial eigenfunctions $\gls{sym-R-1n-i1}(\gls{sym-r})$ and $\gls{sym-R-1m-i2}(\gls{sym-r})$ contain the modified Bessel function of the first kind $\gls{sym-mathrm-I-0}(\gls{sym-chi})$, which diverges for $\gls{sym-chi}\rightarrow \infty$, resulting in numeric overflow.
To prevent overflow, each fluid region must exceed a minimum height $\gls{sym-Delta-z-min}$ that is proportional to the region diameter $\gls{sym-D}$ for both the spar ($\gls{sym-Delta-z-s}, \gls{sym-D-s}$) and the float ($\gls{sym-Delta-z-f}, \gls{sym-D-f}$), with variables shown in \Cref{fig:meem-geom}.
The constant of proportionality between $\gls{sym-Delta-z-min}$ and $\gls{sym-D}$ is a function of the number of harmonics $\gls{sym-N-harmonics}$ used in that region ($\gls{sym-N}$ for spar, $\gls{sym-M}$ for float), as well as the maximum argument to the \texttt{besseli} function in MATLAB before overflow, $\gls{sym-chi-max}$:
\begin{equation}\label{eq:delta-z-min}
    \frac{\gls{sym-Delta-z-min}}{\gls{sym-D}} = \frac{\gls{sym-pi}\gls{sym-N-harmonics}}{2\gls{sym-chi-max}} \approx \frac{\gls{sym-N-harmonics}}{446}.
\end{equation}
By trial, $\gls{sym-chi-max}$ is found to be $\approx 700.5$, close to the theoretical value of \texttt{log (realmax)} $\approx 709.8$ for exponential scaling.
Since $\gls{sym-N-harmonics}=10$ gives adequate convergence for most geometries, this condition is trivially satisfied for nearly all floating bodies of practical relevance, but must still be added as a constraint to prevent the optimizer from exploiting the numerical divergence.
In future work, if high-accuracy solutions are desired for large bodies close to the sea floor, exponentially scaled Bessel functions could be used instead, with symbolic cancellation of the exponential scaler from the eigenfunction numerator and denominator.

Likewise, the vertical eigenfunction $\gls{sym-Z-k-e}$ for $\gls{sym-k-index}=0$ contains the $\cosh$ and $\sinh$ functions, which diverge for large values of $\gls{sym-m-0}\gls{sym-h}$ (high frequencies or deep water).
Since the largest relevant value of $\gls{sym-m-0}\gls{sym-h}$ depends on the site rather than on the \gls{acr-wec} design, it is not possible to add geometric constraints to prevent overflow as it was above.
Therefore, the limit is derived analytically. %\hl{This formula hasn't been numerically checked yet}
\begin{equation}
    \lim_{\gls{sym-m-0}\gls{sym-h}\rightarrow\infty} \gls{sym-Z-0-e} (\gls{sym-z})= \frac{\cosh^2 (\gls{sym-m-0}\gls{sym-h})}{\sqrt{2\gls{sym-m-0}\gls{sym-h}}}\exp\left(1+\frac{\gls{sym-z}}{\gls{sym-h}}\right)
\end{equation}
Plugging this into the first element of the bottom block of the b-vector (position $\gls{sym-N}+2\gls{sym-M}+1$) results in
\begin{equation}
    \lim_{\gls{sym-m-0}\gls{sym-h}\rightarrow\infty}\gls{sym-b-N-2M-1}=\frac{-\gls{sym-a-2}}{\gls{sym-h}-\gls{sym-d-2}} \sqrt{\frac{\gls{sym-h}}{2\gls{sym-m-0}}}\exp(-\gls{sym-d-2}\gls{sym-m-0})
\end{equation}
\begin{figure}[htbp]
    \centering
    \includegraphics[\FigureOptions{aor:figs/from-matlab/asymptotic_b_vector.pdf}]{figs/from-matlab/asymptotic_b_vector.pdf}
    \caption{Asymptotic b-vector for large $\gls{sym-m-0} \gls{sym-h}$}
    \label{fig:meem-b-limit}
\end{figure}

\Cref{fig:meem-b-limit} shows the asymptotic approximation as well as the standard expression for the relevant b-vector entry as a function of nondimensional ratio $\gls{sym-m-0}\gls{sym-h}/\textrm{acosh(realmax})$.
With the standard expression, values for inputs above $\frac{1}{2}$ become \texttt{NaN}, which poses problems in deep water.
With the asymptotic approximation, we are able to avoid \texttt{NaN}s for the zone $\frac{1}{2}<\gls{sym-m-0}\gls{sym-h}/\textrm{acosh(realmax})<\frac{\gls{sym-h}}{\gls{sym-d-2}}$.
Above this threshold, \texttt{NaN}s occur even with the approximation, and the value in the b-vector is set to zero.
This introduces negligible error because in this range, the magnitude of the true value is less than $10^{-300}$ for the nominal \gls{acr-rm3} geometry.
The asymptotic approximation is accurate for values of $\gls{sym-m-0}\gls{sym-h}/\textrm{acosh(realmax)}$ as low as $10^{-2}$,
though it is only applied for values above $\frac{1}{2}$.


The corresponding limit for the vertical coupling integral is:
\begin{equation}
\begin{aligned}
    \lim_{\gls{sym-m-0}\gls{sym-h}\rightarrow\infty}
    \gls{sym-boldsymbol-mathcal-Z}_{\gls{sym-m},\gls{sym-k-generic-index}=0}
    &
    %= \int_{-h}^{-d_2} \vec{Z}_m^{i2 ~T}\vec{Z}_0^{\textit{e}} dz
    = \gls{sym-h}~\frac{
                            \cosh^2 (\gls{sym-m-0}\gls{sym-h})
                        }
                        {
                            \sqrt{2\gls{sym-m-0}\gls{sym-h}}
                        }
            \cdot \frac{
                            -1+(-1)^{\gls{sym-m}}
                            \exp(
                                1-\frac{\gls{sym-d-2}} {\gls{sym-h}}
                                )
                        }
                        {
                            \gls{sym-f-m}
                        }
    \\
    \text{where}~~\gls{sym-f-m}&=
    \begin{cases}
        1,                                   & \gls{sym-m}=0 \\
        \gls{sym-h}^2\gls{sym-lambda-m}^2+1, & \gls{sym-m} \geq 1
    \end{cases}
\end{aligned}
\end{equation}
%\hl{check this formula}

A final numerical subtlety worth discussing is finite precision effects in calculating $\gls{sym-m-k}$.
Bounds of $180^\textrm{o}\cdot[\gls{sym-k-index}-\frac{1}{2}, \gls{sym-k-index}]$
are placed on $\gls{sym-m-k}\gls{sym-h}$ in a root-finding algorithm to ensure the $\gls{sym-k-index}$th root is identified.
Without these bounds, sometimes the \texttt{fsolve} solver would wrongly converge to values where the residual approaches
$+\infty$ from one side and $-\infty$ from the other
(asymptotes of $\tan()$) rather than true zeroes.
Degrees are used instead of radians so asymptotes occur at rational values.

\subsection{Runtime and Computational Cost Scaling}

The runtime of the \gls{acr-meem} method is the time required to find the eigencoefficients, then obtain the hydrodynamic coefficients from eigencoefficients.
First, a nonlinear root-finding algorithm runs $\gls{sym-K}-1$ times to generate the $\gls{sym-m-k}$ inputs used in the A-matrix and b-vector.
Then $3\gls{sym-N}+8\gls{sym-M}+2\gls{sym-K}-11$ Bessel functions must be evaluated for the radial terms of the A-matrix.
 (Note that this would reduce to $2\gls{sym-N}+8\gls{sym-M}+2\gls{sym-K}-10$ evaluations if the definition of the $\gls{sym-R-1n-i1}(\gls{sym-r})$ eigenfunction were modified to use a scale factor based on $\gls{sym-a-1}$ rather than $\gls{sym-a-2}$, although this enhancement is not pursued here to maintain consistency with previous work).
% if I change scaling of i2 from a2 to (a1+a2)/2, it increases by 2* (M-1). So total would be 2N+10M+2K-12. see notebook p51.
The cost of evaluating vertical coupling integrals (\Cref{tab:meem-integrals}) is negligible since they are trigonometric.
Linear solves scale almost cubically with matrix size, so this step scales with $(\gls{sym-N}+2\gls{sym-M}+\gls{sym-K})^3$.
The radial integrals for the c-vector do not require evaluating Bessel functions with any arguments that were not already evaluated for the A-matrix.
For $\gls{sym-N}=\gls{sym-M}=\gls{sym-K}=10$, the simulation of a single frequency averages \resultsAOR[MEEMRuntime] on the hardware described in \Cref{sec:sim-runtime}.
\Cref{fig:runtime-hydro} shows the time breakdown.
Most of the time is spent evaluating Bessel functions, so future code optimization should focus on speeding up Bessel evaluations, such as with lookup tables.
\citet{chau_inertia_2012} proposes using the sparsity pattern to reduce matrix size from N+2M+K to 2M, but this seems low impact since the linear solve only takes a few percent of compute time and the reduced matrix would require an equal number of Bessel evaluations.
On the other hand, matrix size in a boundary element method solver is much larger (meshes may have 1000s of panels) and the linear solve can drive computation cost.
On the same machine, Capytaine boundary element method for the same geometry takes an average of 323 ms for a 710 panel mesh (1\% convergence). %\hl{the 323 ms capy number came from my laptop not from the lab computer}
Thus, \gls{acr-meem} achieves a 10x time reduction over Capytaine.
